\section{Discussion and limitations}
The practical lesson is simple. A statement such as ``the pruned model recovers $x\%$ of its quality'' is incomplete unless the operating point is specified. In this case, the same recovery stage closes $82\%$ of the gap to dense at $10.24$\,s but $44\%$ at $3.84$\,s, and it gives no alignment recovery on the held-out music captions. A single in-domain benchmark at the native duration would describe the favourable part of that behaviour while hiding the rest.

For future compression studies, the evaluation can remain compact while being more informative. The pruned and recovered systems should be paired with common generation noise, recovery should be reported at the training operating point and at deliberately displaced points, and automatic scores should be given interpretable anchors. When domain shift is plausible, at least one held-out prompt battery is useful. These additions test whether recovery behaves like a general restoration or like local specialization without requiring a new pruning method.

The same design is straightforward to reproduce. First establish P and \PFT{} as a matched pair and freeze the generation recipe. Then evaluate both with shared noise at the recovery condition and at one or more controlled departures from it. Compute recovery per prompt before aggregation, because averaging system scores first discards the pairing that reduces sampling variance. Finally, report the paired gain together with a reference gap that gives it scale. The exact durations, domains and scorers will change with the application, but the logic does not. A recovery claim should describe where the gain was measured and how far that evidence extends.

To keep the paper readable, we retain the contrasts needed to support each conclusion and use the figure to show the overall pattern. The companion repository preserves exhaustive sensitivity intervals, scorer-wise values, intermediate operating points, frozen manifests, analysis status and executable provenance for audit and reproduction~\cite{bibbo2026repo}.

The empirical picture suggests a useful way to think about recovery. The recovered checkpoint is not simply P plus a fixed amount of quality. Fine-tuning changes a function whose benefit can depend on the query. In our case, moving away from the recovery conditions along two different axes reveals two different failures of transfer. Shorter requests retain some alignment gain but substantially less of it, whereas the held-out domain shows essentially none. This does not make the recovered checkpoint unhelpful. It makes the scope of the recovery claim more precise.

Our results are consistent with specialization toward the conditions used during recovery fine-tuning, but they do not establish that mechanism. The missing matched control is a dense AudioLDM-M given the same $10^{6}$-step AudioCaps fine-tune. Without it, the duration and domain dependence cannot be attributed uniquely to pruning rather than to fine-tuning in general. We also study only one model family, two pruning severities and one held-out music sub-domain. The music battery differs in caption style as well as content. Primary inference uses CLAP, although the duration pattern reproduces with Human-CLAP and event-level metrics. Author listening is informal and disagrees with CLAP at short duration. The sweep stops at $10.24$\,s, so it cannot distinguish a maximum at the recovery duration from a broader preference for longer clips.

\section{Conclusion}
Recovery fine-tuning is not a single scalar property of a pruned text-to-audio checkpoint. In AudioLDM-M, the recovered gain changes substantially with duration and prompt domain. At $83\%$ pruning it is strongest near the $10.24$\,s in-domain recovery condition, smaller for short generation and absent on held-out hip-hop captions. Evaluating recovery across operating points therefore reveals behaviour that a single benchmark setting can miss, while the paired protocol and repository-level provenance make that conclusion reproducible.

\clearpage
{\setlength{\itemsep}{0pt}\setlength{\parsep}{0pt}
\begin{thebibliography}{99}
\bibitem{fang2023diffpruning} G.~Fang, X.~Ma, and X.~Wang, ``Structural pruning for diffusion
models,'' in \emph{Proc.\ NeurIPS}, 2023.
\bibitem{kim2024bksdm} B.-K.~Kim, H.-K.~Song, T.~Castells \emph{et al.}, ``BK-SDM: A lightweight, fast,
and cheap version of Stable Diffusion,'' in \emph{Proc.\ ECCV}, 2024.
\bibitem{fang2025tinyfusion} G.~Fang, K.~Li, X.~Ma, and X.~Wang, ``TinyFusion: Diffusion transformers
learned shallow,'' in \emph{Proc.\ CVPR}, 2025.
\bibitem{liu2023audioldm} H.~Liu, Z.~Chen, Y.~Yuan \emph{et al.}, ``AudioLDM: Text-to-audio generation with latent diffusion models,'' in
\emph{Proc.\ ICML}, 2023.
\bibitem{singh2026pruning} A.~Singh, Y.~Yuan, Y.~Chen, W.~Wang, and M.~D.~Plumbley, ``Efficient
text-to-audio generation via pruning,'' arXiv:2607.13330, 2026.
\bibitem{lai2021parp} C.-I.~J.~Lai, Y.~Zhang, A.~H.~Liu \emph{et al.}, ``PARP: Prune, adjust and
re-prune for self-supervised speech recognition,'' in \emph{Proc.\ NeurIPS}, 2021.
\bibitem{peng2023dphubert} Y.~Peng, Y.~Sudo, S.~Muhammad, and S.~Watanabe, ``DPHuBERT: Joint
distillation and pruning of self-supervised speech models,'' in \emph{Proc.\ Interspeech}, 2023.
\bibitem{wu2023clap} Y.~Wu, K.~Chen, T.~Zhang \emph{et al.},
``Large-scale contrastive language-audio pretraining with feature fusion and keyword-to-caption
augmentation,'' in \emph{Proc.\ ICASSP}, 2023.
\bibitem{huang2023makeanaudio} R.~Huang, J.~Huang, D.~Yang \emph{et al.}, ``Make-An-Audio:
Text-to-audio generation with prompt-enhanced diffusion models,'' in \emph{Proc.\ ICML}, 2023.
\bibitem{ghosal2023tango} D.~Ghosal, N.~Majumder, A.~Mehrish, and S.~Poria, ``Text-to-audio generation
using instruction guided latent diffusion model,'' in \emph{Proc.\ ACM Int.\ Conf.\ Multimedia (MM)},
2023, pp.~3590--3598.
\bibitem{kilgour2019fad} K.~Kilgour, M.~Zuluaga, D.~Roblek \emph{et al.}, ``Fr\'echet Audio
Distance: A reference-free metric for evaluating music enhancement algorithms,'' in
\emph{Proc.\ Interspeech}, 2019.
\bibitem{gui2024fad} A.~Gui, H.~Gamper, S.~Braun \emph{et al.}, ``Adapting Fr\'echet Audio
Distance for generative music evaluation,'' in \emph{Proc.\ ICASSP}, 2024.
\bibitem{chung2025kad} Y.~Chung, P.~Eu, J.~Lee \emph{et al.}, ``KAD: No more FAD!
An effective and efficient evaluation metric for audio generation,'' arXiv:2502.15602, 2025.
\bibitem{kong2020panns} Q.~Kong, Y.~Cao, T.~Iqbal \emph{et al.}, ``PANNs:
Large-scale pretrained audio neural networks for audio pattern recognition,'' \emph{IEEE/ACM
Trans.\ Audio, Speech, Lang.\ Process.}, vol.~28, pp.~2880--2894, 2020.
\bibitem{takano2025humanclap} T.~Takano, Y.~Okamoto, Y.~Kanamori \emph{et al.}, ``Human-CLAP:
Human-perception-based contrastive language-audio pretraining,'' in \emph{Proc.\ APSIPA ASC}, 2025,
pp.~131--136.
\bibitem{li2026finelap} X.~Li, X.~Xu, Z.~Ma \emph{et al.}, ``FineLAP: Taming
heterogeneous supervision for fine-grained language-audio pretraining,'' in \emph{Proc.\ ACL}, 2026,
pp.~10393--10408.
\bibitem{liu2024audioldm2} H.~Liu, Y.~Yuan, X.~Liu \emph{et al.}, ``AudioLDM 2: Learning holistic audio generation with self-supervised
pretraining,'' \emph{IEEE/ACM Trans.\ Audio, Speech, Lang.\ Process.}, vol.~32, pp.~2871--2883, 2024.
\bibitem{kumar2022finetuning} A.~Kumar, A.~Raghunathan, R.~Jones \emph{et al.}, ``Fine-tuning
can distort pretrained features and underperform out-of-distribution,'' in \emph{Proc.\ ICLR}, 2022.
\bibitem{kim2019audiocaps} C.~D.~Kim, B.~Kim, H.~Lee, and G.~Kim, ``AudioCaps: Generating captions
for audios in the wild,'' in \emph{Proc.\ NAACL-HLT}, 2019.
\bibitem{agostinelli2023musiclm} A.~Agostinelli, T.~I.~Denk, Z.~Borsos \emph{et al.}, ``MusicLM:
Generating music from text,'' arXiv:2301.11325, 2023.
\bibitem{song2021ddim} J.~Song, C.~Meng, and S.~Ermon, ``Denoising diffusion implicit models,'' in
\emph{Proc.\ ICLR}, 2021.
\bibitem{wang2026ttabench} H.~Wang, C.~Liu, J.~Chen \emph{et al.}, ``TTA-Bench: A comprehensive benchmark for evaluating text-to-audio models,'' in \emph{Proc.\ AAAI}, 2026, pp.~33512--33520.
\bibitem{bibbo2026repo} G.~Bibb\'o, A.~Singh, and M.~D.~Plumbley, ``Companion repository for Recovery Fine-Tuning Recovers Where It Was Trained,'' 2026. [Online]. Available: \url{https://github.com/gbibbo/audioldm-modality-swap-pruning/blob/main/PAPER_COMPANION.md}
\end{thebibliography}}

\end{document}
